\documentclass[10pt,a4paper]{amsart}
\usepackage[margin=19mm]{geometry}
\usepackage{amsmath,amssymb,mathtools}
\usepackage[scr=rsfs,cal=euler]{mathalfa}
\usepackage{xfrac}
\usepackage[unicode,hidelinks,bookmarks=false]{hyperref}
\newcommand{\ie}{i.e.\ }
\newcommand{\cf}{cf.\ }
\newcommand{\resp}{respectively\ }
\newcommand{\Subsection}{\subsection}
\providecommand{\nobreakdash}{\nobreak\hbox{-}}
\setlength{\emergencystretch}{2em}
\multlinegap0pt
\usepackage{mathtools}
\usepackage[normalem]{ulem}
\usepackage{tikz}
\usepackage{dsfont}
\usepackage{amssymb}
\usepackage{braket}
\usepackage{cancel}
\usepackage{float}
\newtheorem{lem}{Lemma}
\title{The completion of the set of Lagrangians \\ and applications to dynamics \\ \Large Based on lectures by C. Viterbo}
\author{Olga Bernardi, Francesco Morabito}
\date{Source: arXiv:2603.09396v1 (CC BY 4.0)}
\begin{document}
\maketitle

\noindent\emph{Source and licence.} The completion of the set of Lagrangians \\ and applications to dynamics \\ \Large Based on lectures by C. Viterbo. Version arXiv:2603.09396v1 (CC BY 4.0), distributed by arXiv under the Creative Commons Attribution 4.0 International licence, http://creativecommons.org/licenses/by/4.0/, as recorded in the arXiv licence metadata for this exact version and shown on the abstract page https://arxiv.org/abs/2603.09396v1. Full licence notice: \texttt{licenses/arxiv-2603.09396-CC-BY-4.0.txt}.\par\smallskip

\noindent\textit{Editorial scope.} Counted unit: the article's lem lem (source0.tex lines 487--489). The article's complete original proof of it is delivered with it (source0.tex lines 490--495, one \texttt{proof} environment of 6 lines). No mathematics is written, repaired or re-derived: the statement and the proof are the article's own lines, in the article's own notation, and no retained line is substituted or re-flowed.  No further source-local context is needed: every cross-reference of every retained line resolves inside the two delivered spans.  Nothing was omitted from any retained span, so the package carries no omission record.  No retained line contains a citation, so the wrapper carries no bibliography and no bibliography entry.  No retained line contains a figure environment, an image-inclusion command or drawing code, so no image is delivered and the package declares no asset dependency.  Editorial formatting only, in addition: an AMS \texttt{amsart} preamble that restates the article's own declarations and macros used by the retained lines, namely the environments \texttt{lem} on separate counters, together with the standard packages those lines need.  The retained environments print in this document as Lemma 1 in order of appearance, whereas the article prints the same items with the numbering of its own full document.  The complete native source of the article as distributed in the arXiv e-print for this exact version is bundled unchanged as \texttt{sources/196056/source0.tex}.
\begin{lem}
Let $S\in C^{\infty}(N \times \mathbb{R}^k,\mathbb{R})$, $(q,\xi) \mapsto S(q;\xi)$ be a GFQI. Then $S$ is PS.
\end{lem}

\begin{proof}
Let $(q_n, \xi_n)_{n \ge 1}$ be a sequence satisfying
$$\vert{dS(q_n,\xi_n)}\vert \to 0 \qquad {and} \qquad S(q_n,\xi_n) \text{ bounded}\, .$$
If the sequence $(\xi_n)_{n \ge 1}$ is contained in a compact set, then $(q_n, \xi_n)_{n\ge 1}$ admits 
a converging subsequence (because $N$ is compact). Let verify that this is the unique possible case. Remember that $S(q_n,\xi_n) = \xi_n^T Q \xi_n$ for $|\xi_n| \ge R$. It there were only finite terms of the sequence $(\xi_n)_{n \ge 1}$ in every compact set, it would follow that $\lim_{n \to +\infty} |\xi_n| = +\infty$. But this is not possible because, since $\partial_{\xi}S(q,\xi_n) = 2 Q \xi_n$ (for $|\xi_n| \ge R$) and $Q$ is non degenerate, $\vert{dS(q_n,\xi_n)}\vert$ would tend to $+\infty$ and not to $0$.
\end{proof}

\end{document}
